\section{Closed-loop protocol with drive tracking}
\label{sec:feedback}

The control objective is to find a flux offset $V$ such that
\begin{equation}
 r(V)=\omega_q(V)-\omega_{\mathrm{tar}}=0.
 \label{eq:root-problem}
\end{equation}
At iteration $n$, the transient discriminator returns
$\widehat\omega_{q,n}=\omega_{d,n}+\widehat\Delta_n$.  The estimated feedback
residual is $\widehat r_n=\widehat\omega_{q,n}-\omega_{\mathrm{tar}}$.  A
damped gradient update changes the flux bias according to
\begin{equation}
 V_{n+1}=\operatorname{clip}\left(
 V_n-\eta\frac{\widehat r_n}{\widehat S_n},
 V_{\min},V_{\max}\right),
 \label{eq:bias-update}
\end{equation}
where $\eta$ is a damping factor and $\widehat S_n$ is a local slope.  After
two measurements, a model-free secant estimate is
\begin{equation}
 \widehat S_n=\frac{\widehat\omega_{q,n}-\widehat\omega_{q,n-1}}
 {V_n-V_{n-1}}.
 \label{eq:secant-slope}
\end{equation}

The measurement reference is updated separately from the controlled bias.  We
predict the qubit frequency at the next point using
\begin{equation}
 \omega_{d,n+1}=\widehat\omega_{q,n}
 +\widehat S_n(V_{n+1}-V_n).
 \label{eq:drive-tracking}
\end{equation}
Equation~\eqref{eq:drive-tracking} serves only to keep
$|\omega_q(V_{n+1})-\omega_{d,n+1}|$ within the local discriminator window.
It does not change the optimization target in Eq.~\eqref{eq:root-problem}.  An
initial frequency estimate seeds the first drive, and a bounded blind step
provides the second point needed for Eq.~\eqref{eq:secant-slope}.  Subsequent
updates use the measured secant slope, so the procedure does not require an
analytic frequency--flux model.

The complete protocol therefore has two coupled but distinct loops: the flux
update reduces the target error, whereas the drive update reduces the
measurement detuning.  A range guard can trigger a wide-capture Ramsey
measurement if the local discriminator leaves its validated interval.  The
numerical experiment below isolates the local tracking phase and starts from a
frequency seed sufficient to keep the transient probe on its connected branch.
